%%%%%%%%%%%%%%%%%%%%========== ハイパーリンク ==========%%%%%%%%%%%%%%%%%%%%
% hyperrefはほかのパッケージより後に，cleverefはhyperrefより後に読み込む．
\usepackage[unicode, hidelinks]{hyperref}
\hypersetup{%
	pdftitle = {\booktitle},
	pdfauthor = {\bookauthor},
	bookmarksnumbered = true,
	bookmarksopen = true,
	bookmarkstype = toc,
	pdfborder = {0 0 0},
	colorlinks = true,
}
\usepackage{footnotebackref}

%%%%%%%%%%%%%%%%%%%%========== 相互参照 ==========%%%%%%%%%%%%%%%%%%%%
\usepackage[nameinlink]{cleveref}

% \crefと\Crefに同じ名前を設定する．単数形と複数形も区別しない．
\NewDocumentCommand{\setcrefname}{mm}{%
	\crefname{#1}{#2}{#2}%
	\Crefname{#1}{#2}{#2}%
}

% 定理環境．theorems.texに環境を追加したら，ここにも追加する．
\setcrefname{Def}{定義}
\setcrefname{Thm}{定理}
\setcrefname{Prop}{命題}
\setcrefname{Lemma}{補題}
\setcrefname{Corollary}{系}
\setcrefname{Ex}{例}
\setcrefname{Note}{注意}
\setcrefname{Que}{演習}

% 図表・数式・ページ・箇条書き
\setcrefname{figure}{図}
\setcrefname{table}{表}
\setcrefname{equation}{式}
\setcrefname{page}{p.}
\setcrefname{enumi}{}

% 章・節・小節
\crefformat{chapter}{#2第#1章#3}
\Crefformat{chapter}{#2第#1章#3}
\crefrangeformat{chapter}{#3第#1#4章から#5第#2#6章}
\Crefrangeformat{chapter}{#3第#1#4章から#5第#2#6章}
\setcrefname{section}{\presectionname}
\setcrefname{subsection}{\presubsectionname}
\crefrangeformat{section}{#3\presectionname#1#4から#5\presectionname#2#6}
\Crefrangeformat{section}{#3\presectionname#1#4から#5\presectionname#2#6}
\crefrangeformat{subsection}{#3\presubsectionname#1#4から#5\presubsectionname#2#6}
\Crefrangeformat{subsection}{#3\presubsectionname#1#4から#5\presubsectionname#2#6}

%%%%%%%%%%%%%%%%%%%%========== 参考文献 ==========%%%%%%%%%%%%%%%%%%%%
\usepackage[%
	backend = biber,%
	url = false,%
	doi = false,%
	eprint = false,%
	isbn = false,%
	sorting = none,%
	style = numeric-comp%
]{biblatex}
\DeclareFieldFormat*{title}{\mkbibquote{#1}}
\DeclareDelimFormat{labelnamepunct}{\addcomma\addspace}
\DeclareFieldFormat{url}{\url{#1}}
\addbibresource{reference/book.bib}
